\chapter{Totale Beschränktheit und Kompaktheit}
\hypertarget{compactness.main}{}

Die diskrete Metrik auf \(\mathbb N\) zeigt, dass selbst Vollständigkeit
und Beschränktheit zusammen keine Folgenkompaktheit erzwingen. Die fehlende
Bedingung betrifft die Verteilung der Punkte bei beliebig kleinen Radien:
Für jede Genauigkeit sollen endlich viele Kugeln den ganzen Raum erfassen.
Dies führt zur totalen Beschränktheit und anschließend zum
Überdeckungsbegriff der Kompaktheit.

Die \CompactnessReadingLink{} entwickelt den Zusammenhang mit Beispielen
und einem zusammenhängenden Beweis. Die \CompactnessProofsLink{} enthalten
die vollständigen Einzelbeweise einschließlich der Auswahl- und
Rekursionsschritte. Die Aussagen und Definitionen bleiben in diesem Band.

\section{Endliche Netze und totale Beschränktheit}

\FormulaDefDeltaK[Endliches Epsilon-Netz]{%
  \begin{aligned}[t]
  &\operatorname{Metrik}_{X}(d)\dsep K\subseteq X\dsep\varepsilon>0
    \vdash\\[-2pt]
  &\quad\operatorname{Netz}_{d}(F,K,\varepsilon)\coloneqq\\[-2pt]
  &\qquad F\subseteq K\land\Finite{F}\land
    \forall x\in K\,\exists c\in F\;(d(x,c)<\varepsilon).
  \end{aligned}
}{MetricFiniteEpsilonNetDef}{%
  \DeltaRow{Mengen}{X\dsep K\dsep F\dsep x\dsep c\dsep\varepsilon}
  \DeltaRow{Funktionen}{d}
  \DeltaRow{Neues Prädikat}{\operatorname{Netz}_{d}}
}

Die Zentren gehören zu \(K\). Gleichwertig ist
\(K\subseteq\bigcup_{c\in F}B_\varepsilon^d(c)\).
Für \(K=\varnothing\) ist \(F=\varnothing\) bei jedem positiven Radius
ein Netz. Da die Abstände innerhalb von \(K\) unverändert bleiben,
stimmt der Netzbegriff für \(d\) mit demjenigen für \(d_K\) überein.

\FormulaDefDeltaK[Total beschränkte Teilmenge]{%
  \begin{aligned}[t]
  &\operatorname{Metrik}_{X}(d)\dsep K\subseteq X\vdash\\[-2pt]
  &\quad\operatorname{Totalbeschr}_{d}(K)\coloneqq
    \forall\varepsilon\in\mathbb R\,
    \bigl(0<\varepsilon\rightarrow
      \exists F\,\operatorname{Netz}_{d}(F,K,\varepsilon)\bigr).
  \end{aligned}
}{MetricTotalBoundednessDef}{%
  \DeltaRow{Mengen}{X\dsep K\dsep F\dsep\varepsilon}
  \DeltaRow{Funktionen}{d}
  \DeltaRow{Neues Prädikat}{\operatorname{Totalbeschr}_{d}}
}

Jede endliche Teilmenge ist total beschränkt: Man verwendet sie selbst als
Netz. Bei \((\mathbb N,\delta_{\mathbb N})\) sind dagegen die Kugeln vom
Radius \(1/2\) Einermengen; endlich viele davon überdecken \(\mathbb N\)
nicht. Dieses Beispiel ist somit beschränkt, aber nicht total beschränkt.

\FormulaThmDeltaK[Totale Beschränktheit impliziert Beschränktheit]{%
  \operatorname{Metrik}_{X}(d)\dsep K\subseteq X\dsep
  \operatorname{Totalbeschr}_{d}(K)
  \vdash\operatorname{Beschr}_{d}(K)
}{MetricTotalBoundedSetBounded}{%
  \DeltaRow{Mengen}{X\dsep K}
  \DeltaRow{Funktionen}{d}
}
\noindent\emph{Beweis:} \CompactnessProofLink{MetricTotalBoundedSetBounded}.

\section{Cauchy-Teilfolgen}

Totale Beschränktheit garantiert noch keinen Grenzwert im Raum. Sie
garantiert zunächst, dass aus jeder Folge eine Teilfolge ausgewählt werden
kann, deren späte Glieder beliebig nahe beieinander liegen.

\FormulaThmDeltaK[Teilfolgenkriterium für totale Beschränktheit]{%
  \begin{aligned}[t]
  &\operatorname{Metrik}_{X}(d)\dsep K\subseteq X\vdash\\[-2pt]
  &\quad\operatorname{Totalbeschr}_{d}(K)\leftrightarrow\\[-2pt]
  &\qquad\forall a\colon\mathbb N\to K\,
    \exists\varphi\colon\mathbb N\to\mathbb N\,
    \Bigl(\forall n\in\mathbb N\;(\varphi(n)<\varphi(n+1))\\[-2pt]
  &\hspace{56mm}\land\operatorname{Cauchy}_{d_K}(a\circ\varphi)\Bigr).
  \end{aligned}
}{MetricTotalBoundedCauchySubsequenceCriterion}{%
  \DeltaRow{Mengen}{X\dsep K\dsep n}
  \DeltaRow{Funktionen}{d\dsep d_K\dsep a\dsep\varphi}
}
\noindent\emph{Beweis:} \CompactnessProofLink{MetricTotalBoundedCauchySubsequenceCriterion}.

\section{Vollständigkeit und Folgenkompaktheit}

Der folgende Satz verbindet die beiden Aufgaben: Totale Beschränktheit
liefert eine Cauchy-Teilfolge, Vollständigkeit ihren Grenzwert in \(K\).

\FormulaThmDeltaK[Charakterisierung der Folgenkompaktheit]{%
  \begin{aligned}[t]
  &\operatorname{Metrik}_{X}(d)\dsep K\subseteq X\vdash\\[-2pt]
  &\quad\operatorname{Folgenkompakt}_{d}(K)
    \leftrightarrow
    \bigl(\operatorname{Vollst}(K,d_K)\land
          \operatorname{Totalbeschr}_{d}(K)\bigr).
  \end{aligned}
}{MetricSequentialCompactnessCharacterization}{%
  \DeltaRow{Mengen}{X\dsep K}
  \DeltaRow{Funktionen}{d\dsep d_K}
}
\noindent\emph{Beweis:} \CompactnessProofLink{MetricSequentialCompactnessCharacterization}.

Die Vollständigkeit betrifft den Teilraum \((K,d_K)\), nicht notwendig den
umgebenden Raum \(X\). Die beiden Bedingungen erfüllen unterschiedliche
Aufgaben: Das offene reelle Intervall \((0,1)\) ist total beschränkt,
aber unvollständig; der unendliche diskrete Raum ist vollständig, aber
nicht total beschränkt. Die \CompactnessReadingLink{} führt beide Beispiele aus.

\section{Kompaktheit durch offene Überdeckungen}

Eine offene Überdeckung von \(K\) ist eine Familie \(\mathcal U\) von
Teilmengen von \(K\), die bezüglich der Teilraummetrik \(d_K\) offen sind
und zusammen ganz \(K\) ergeben. Damit ist die Definition unabhängig
davon, in welchen größeren metrischen Raum \(K\) eingebettet wird.

\FormulaDefDeltaK[Kompakte Teilmenge]{%
  \begin{aligned}[t]
  &\operatorname{Metrik}_{X}(d)\dsep K\subseteq X\vdash\\[-2pt]
  &\quad\operatorname{Kompakt}_{d}(K)\coloneqq
    \forall\mathcal U\subseteq\mathcal P(K)\,\Bigl(\\[-2pt]
  &\qquad\bigl(\bigcup\mathcal U=K\land
      \forall U\in\mathcal U\;\operatorname{Offen}_{d_K}(U)\bigr)
      \rightarrow\\[-2pt]
  &\qquad\exists\mathcal V\subseteq\mathcal U\,
      \bigl(\Finite{\mathcal V}\land\bigcup\mathcal V=K\bigr)\Bigr).
  \end{aligned}
}{MetricOpenCoverCompactnessDef}{%
  \DeltaRow{Mengen}{X\dsep K\dsep\mathcal U\dsep\mathcal V\dsep U}
  \DeltaRow{Funktionen}{d\dsep d_K}
  \DeltaRow{Neues Prädikat}{\operatorname{Kompakt}_{d}}
}

Die leere Teilüberdeckung ist erlaubt. Insbesondere ist die leere Menge
kompakt. Für allgemeine metrische Räume stimmen nun die drei folgenden
Beschreibungen überein.

\Needspace{14\baselineskip}
\FormulaThmDeltaK[Drei Charakterisierungen metrischer Kompaktheit]{%
  \begin{aligned}[t]
  &\operatorname{Metrik}_{X}(d)\dsep K\subseteq X\vdash\\[-2pt]
  &\quad\bigl(\operatorname{Kompakt}_{d}(K)
      \leftrightarrow\operatorname{Folgenkompakt}_{d}(K)\bigr)\\[-2pt]
  &\quad\land\bigl(\operatorname{Kompakt}_{d}(K)
      \leftrightarrow
      (\operatorname{Vollst}(K,d_K)\land
       \operatorname{Totalbeschr}_{d}(K))\bigr).
  \end{aligned}
}{MetricCompactnessEquivalences}{%
  \DeltaRow{Mengen}{X\dsep K}
  \DeltaRow{Funktionen}{d\dsep d_K}
}
\noindent\emph{Beweis:} \CompactnessProofLink{MetricCompactnessEquivalences}.

Damit sind insbesondere die abgeschlossenen beschränkten reellen Intervalle
aus \FormulaRefAuto{RealClosedBoundedIntervalSequentiallyCompact}[thm]
auch im Überdeckungssinn kompakt. Die zuvor bewiesenen Aussagen über
folgenkompakte Mengen und ihre stetigen Bilder gelten folglich ebenfalls
für kompakte Mengen. Normen oder Vektorraumoperationen werden für diese
Zusammenhänge nicht benötigt.
